% TEMPLATE-MILC-v3.tex - plantilla maestra para todas las guias MILC v3
% Ver scripts/generadores/build-guia-milc.py para la lista completa de
% claves esperadas. El script valida que no queden placeholders sin
% reemplazar antes de escribir el archivo final.

\documentclass[11pt,letterpaper]{article}
\usepackage[margin=1.75cm]{geometry}
\usepackage[table]{xcolor}
\usepackage{fontspec}
\usepackage[spanish]{babel}
\usepackage{tikz}
\usepackage[most]{tcolorbox}
\usepackage{tabularx}
\usepackage{array}
\usepackage{fancyhdr}
\usepackage{titlesec}
\usepackage{ragged2e}
\usepackage{enumitem}
\usepackage{microtype}

\setmainfont{Helvetica Neue}
\setsansfont{Helvetica Neue}
\setlength{\parindent}{0pt}
\setlength{\parskip}{6pt}
\hyphenpenalty=200
\exhyphenpenalty=200
\pretolerance=400
\tolerance=2000
\setlength{\emergencystretch}{1em}
\renewcommand{\arraystretch}{1.25}
\setlist{leftmargin=5mm,itemsep=1mm,topsep=1mm}
\newcolumntype{Y}{>{\RaggedRight\arraybackslash}X}

\definecolor{milcMagenta}{HTML}{E6007E}
\definecolor{milcVino}{HTML}{5A0038}
\definecolor{milcTurquesa}{HTML}{0093A5}
\definecolor{milcVerde}{HTML}{58A923}
\definecolor{milcMostaza}{HTML}{E5B400}
\definecolor{milcOcre}{HTML}{B86600}
\definecolor{milcNegro}{HTML}{241816}
\definecolor{milcGris}{HTML}{F7F7F4}
\definecolor{milcLinea}{HTML}{E8E2DD}
\definecolor{milcGradePrimary}{HTML}{FF2D87}
\definecolor{milcGradeDark}{HTML}{9F1B56}
\definecolor{milcGradeSoft}{HTML}{FFE0EE}
\definecolor{milcGradeAccent}{HTML}{CFE7FF}
\definecolor{milcGradeText}{HTML}{FFFFFF}
\color{milcNegro}

\pagestyle{fancy}
\fancyhf{}
\lhead{\small Tecnología e Informática · Grado 10}
\rhead{\small Guía 18 · MILC v3}
\cfoot{\small\thepage}
\renewcommand{\headrulewidth}{0pt}

\newtcolorbox{titlebox}[1]{
  enhanced,colback=#1,colframe=#1,arc=7mm,boxrule=0pt,
  left=7mm,right=7mm,top=3mm,bottom=3mm,
  before skip=8pt,after skip=8pt,
  fontupper=\sffamily\bfseries\Large\color{white}
}

\newtcolorbox{softbox}[3]{
  enhanced,colback=#2,colframe=#1,borderline west={4pt}{0pt}{#1},
  arc=2mm,boxrule=.4pt,left=4mm,right=4mm,top=3mm,bottom=3mm,
  before skip=6pt,after skip=8pt,title={#3},coltitle=white,
  colbacktitle=#1,fonttitle=\sffamily\bfseries,fontupper=\RaggedRight,
  attach boxed title to top left={xshift=3mm,yshift=-2mm},
  boxed title style={arc=2mm, boxrule=0pt}
}

\newtcolorbox{drawbox}[2]{
  enhanced,colback=white,colframe=#1,arc=4mm,boxrule=1pt,
  left=5mm,right=5mm,top=4mm,bottom=4mm,before skip=8pt,after skip=8pt,
  title={#2},coltitle=white,colbacktitle=#1,fonttitle=\sffamily\bfseries,
  fontupper=\RaggedRight,
  attach boxed title to top left={xshift=4mm,yshift=-2mm},
  boxed title style={arc=3mm, boxrule=0pt}
}

\newtcolorbox{infoband}[1]{
  enhanced,colback=#1!8,colframe=#1!50,arc=2mm,boxrule=.5pt,
  left=4mm,right=4mm,top=2mm,bottom=2mm,before skip=4pt,after skip=4pt,
  fontupper=\small\RaggedRight
}

% Puente narrativo entre secciones (contrato editorial Conéctate).
% Da continuidad: "Acabas de X. Ahora vas a Y porque Z."
% Visualmente: banda izquierda en color del grado, texto en itálica pequeña.
\newtcolorbox{puentebox}{
  enhanced,colback=milcGradeSoft,colframe=milcGradeDark,
  borderline west={3pt}{0pt}{milcGradeDark},
  arc=0mm,boxrule=0pt,left=5mm,right=4mm,top=2.5mm,bottom=2.5mm,
  before skip=4pt,after skip=8pt,
  fontupper=\itshape\small\color{milcNegro}\RaggedRight
}
\newcommand{\puente}[1]{%
  \begin{puentebox}%
    \textcolor{milcGradeDark}{\textbf{\textrightarrow}}\hspace{2mm}#1%
  \end{puentebox}%
}

\newcommand{\linead}[1]{\noindent\color{milcLinea}\rule{#1}{0.5pt}\color{milcNegro}}
\newcommand{\respuesta}[1]{\par\vspace{2mm}\foreach \n in {1,...,#1}{\noindent\color{milcLinea}\rule{\linewidth}{0.5pt}\par\vspace{5mm}}\color{milcNegro}}
\newcommand{\checkbox}{\textcolor{milcGradeDark}{\fbox{\phantom{x}}}\hspace{5pt}}

\newcommand{\phasecard}[4]{
\begin{tcolorbox}[enhanced,colback=#3,colframe=#2,arc=3mm,boxrule=.5pt,height=3.05cm,valign=center,halign=center,left=1.5mm,right=1.5mm,top=2mm,bottom=2mm]
{\sffamily\bfseries\fontsize{22}{24}\selectfont\textcolor{#2}{#1}}\par
{\sffamily\bfseries\small #4}
\end{tcolorbox}
}

\begin{document}
\thispagestyle{empty}
\begin{tikzpicture}[remember picture,overlay]
  \fill[white] (current page.south west) rectangle (current page.north east);
  \path[fill=milcGradePrimary,rounded corners=16mm]
    ([xshift=.55cm,yshift=-.55cm]current page.north west) rectangle
    ([xshift=-.55cm,yshift=3.65cm]current page.south east);

  \node[anchor=north west,text=milcGradeText] at ([xshift=2.0cm,yshift=-1.85cm]current page.north west)
    {\sffamily\bfseries\fontsize{14}{16}\selectfont GRADO};
  \node[anchor=north west,text=milcGradeText,opacity=.82] at ([xshift=2.0cm,yshift=-2.75cm]current page.north west)
    {\sffamily\bfseries\fontsize{9}{11}\selectfont SERIE GUÍAS · TECNOLOGÍA E INFORMÁTICA};

  \begin{scope}[shift={([xshift=-3.05cm,yshift=-2.55cm]current page.north east)}]
    \fill[white,opacity=.80] (-.62,-.52) rectangle (.62,.52);
    \draw[milcGradeText,line width=1pt] (-.62,-.52) rectangle (.62,.52);
    \fill[milcGradeDark] (-.42,-.38) rectangle (-.22,.06);
    \fill[milcGradeAccent] (-.10,-.38) rectangle (.10,.24);
    \fill[milcGradePrimary!60!black] (.22,-.38) rectangle (.42,.38);
  \end{scope}

  \node[anchor=west,text=milcGradeText] at ([xshift=1.9cm,yshift=-12.85cm]current page.north west)
    {\sffamily\bfseries\fontsize{124}{124}\selectfont 10};
  \draw[milcGradeText,line width=1.7pt]
    ([xshift=4.52cm,yshift=-11.68cm]current page.north west) circle (.13cm);

  \node[anchor=west,text=milcGradeText,text width=8.7cm,align=left] at ([xshift=10.35cm,yshift=-11.6cm]current page.north west)
    {
     {\sffamily\bfseries\fontsize{19}{22}\selectfont Correo profesional ---\\con IA como asistente,\\no como reemplazo}\\[4mm]
     {\sffamily\bfseries\fontsize{23}{26}\selectfont Guía 18-10-TIC}\\[2mm]
     {\sffamily\fontsize{13}{16}\selectfont Período 2 · Informes técnicos y comunicación profesional con IA}\\[5mm]
     {\sffamily\bfseries\fontsize{11}{13}\selectfont Institución Educativa Sor María Juliana}\\[1mm]
     {\sffamily\fontsize{10}{12}\selectfont Autor: Dr. Álvaro Cárdenas Orozco}
    };

  \path[fill=milcGradeSoft,rounded corners=7mm]
    ([xshift=1.35cm,yshift=.85cm]current page.south west) rectangle
    ([xshift=-1.35cm,yshift=4.25cm]current page.south east);
  \draw[milcGradeDark,line width=.65pt,rounded corners=7mm]
    ([xshift=1.35cm,yshift=.85cm]current page.south west) rectangle
    ([xshift=-1.35cm,yshift=4.25cm]current page.south east);
  \node[anchor=center,text=milcNegro,text width=17.0cm,align=center] at ([yshift=2.55cm]current page.south)
    {\sffamily\fontsize{10}{12}\selectfont Metodología MILC · Apertura ancestral · Pensamiento computacional · Triángulo Dussel-Estoicismo-Floridi\\
    \textcolor{milcGradeDark}{\bfseries Producto:} 3 correos profesionales escritos en colaboración con IA (uno solicitando información a una institución o persona, uno respondiendo a un cliente o usuario, uno declinando una propuesta con respeto) + bitácora detallada con 4 columnas por cada correo (a) el prompt usado, (b) la versión generada por la IA, (c) la versión final tras tu edición humana, (d) por qué cambiaste lo que cambiaste + ficha de 5 elementos del correo profesional ideal para uso futuro};
\end{tikzpicture}
\mbox{}
\newpage

\section*{Apertura: Correo profesional con IA --- asistente, no reemplazo}

\begin{softbox}{milcGradeDark}{milcGradeSoft}{Saber ancestral}
En los pueblos del Valle del Cauca y del Quindío cafetero, antes del correo electrónico, hubo un oficio crucial para la comunicación a distancia: \textbf{el telegrafista del pueblo}. Sentado en la estación frente al aparato Morse, el telegrafista enviaba mensajes a través del país. Su negocio tenía regla simple: \textbf{cobraba por palabra}. Cada palabra costaba pesos al cliente. Por eso, el oficio del telegrafista entrenaba al cliente en algo que muchos modernos olvidan: \textbf{decir lo justo}. Si una señora quería avisar a su hija en Bogotá que iba en bus, no enviaba: \emph{``Querida hija, te escribo para informarte que el día de mañana saldré en el autobús de las 6...''}. Enviaba: \emph{``Salgo bus seis mañana. Llego cinco tarde. Mamá''}. \textbf{Siete palabras}. Toda la información necesaria. Cero desperdicio. La sabiduría del telegrafista era precisa: \textbf{máximo respeto al lector, mínimo desperdicio de palabras}. El correo electrónico profesional hereda esa disciplina: la pantalla del lector es la billetera del telegrafista. Cada palabra extra cuesta atención. La IA ayuda a redactar pero \textbf{tú decides el tono}.
\end{softbox}

\begin{softbox}{milcTurquesa}{milcGris}{Saber contemporáneo}
El \textbf{correo profesional electrónico} es el formato de comunicación escrita más usado del mundo laboral. Sus 5 elementos obligatorios son: (1) \textbf{Asunto claro}: en 5-7 palabras, qué espera el destinatario al abrir. \emph{``Propuesta de mejora del recreo''} mejor que \emph{``Hola''}. (2) \textbf{Saludo apropiado}: \emph{``Estimado/a [nombre]''} para formal, \emph{``Hola [nombre]''} para semi-formal. Evita \emph{``A quien corresponda''} salvo cuando no se conoce destinatario. (3) \textbf{Cuerpo conciso}: 3-5 párrafos cortos. Primer párrafo: motivo del correo. Párrafos del medio: detalles. Último párrafo: pedido o acción esperada. (4) \textbf{Cierre cordial}: \emph{``Cordialmente''}, \emph{``Quedo atento/a''}, \emph{``Gracias por su atención''}. (5) \textbf{Firma con datos}: nombre completo, cargo o rol (estudiante grado 10°), institución (IE Sor María Juliana), contacto si aplica. La regla profesional: \textbf{si el correo pasa de 200 palabras, conviene revisar si todo el contenido es necesario}. La IA puede ayudar a redactar el borrador inicial, pero el editor humano decide tono, pedido y matices culturales.
\end{softbox}

\begin{softbox}{milcMostaza}{milcGris}{Pregunta-puente}
\textit{¿Qué sabía la señora del barrio al enviar telegrama de 7 palabras con toda la información necesaria, que el novato olvida cuando manda correo de 500 palabras sin pedido claro? ¿Y por qué la IA puede generar correos correctos pero culturalmente fríos para audiencias colombianas?}
\end{softbox}

\begin{softbox}{milcVerde}{milcGris}{Saber y saber hacer}
Al terminar podrás: (1) \textbf{identificar} los 5 elementos obligatorios del correo profesional y qué tono usar según la situación; (2) \textbf{aplicar} la IA como asistente para redactar borradores iniciales sin renunciar al tono propio; (3) \textbf{analizar} las diferencias entre versión IA y versión final con tu voz; (4) \textbf{evaluar} qué tan profesional queda el correo final con criterio cultural colombiano.
\end{softbox}

\small
\begin{tabularx}{\textwidth}{>{\bfseries}p{3.0cm}Y}
\rowcolor{milcGradePrimary!12}Autor & Dr. Álvaro Cárdenas Orozco\\
DBA & Comunicación profesional con asistencia de IA y herramientas ofimáticas gratuitas (MEN, T\&I)\\
Referentes & Markdown como estándar abierto · Google Workspace · Estoicismo (claridad)\\
Producto final & 3 correos profesionales escritos en colaboración con IA (uno solicitando información a una institución o persona, uno respondiendo a un cliente o usuario, uno declinando una propuesta con respeto) + bitácora detallada con 4 columnas por cada correo (a) el prompt usado, (b) la versión generada por la IA, (c) la versión final tras tu edición humana, (d) por qué cambiaste lo que cambiaste + ficha de 5 elementos del correo profesional ideal para uso futuro\\
\end{tabularx}
\normalsize

\puente{Antes de empezar, mira el viaje completo. En la sesión 7 escribiste informe comercial. Hoy aprendes a \emph{enviar} ese informe profesionalmente o cualquier otro mensaje del mundo laboral. Las próximas sesiones (mini-proyecto, sustentación) consolidarán todo el periodo.}

\begin{titlebox}{milcGradeDark}
Ruta MILC: así vamos a aprender
\end{titlebox}
\begin{tabularx}{\textwidth}{>{\centering\arraybackslash}X>{\centering\arraybackslash}X>{\centering\arraybackslash}X>{\centering\arraybackslash}X}
\phasecard{1}{milcVerde}{milcGris}{Escuta\\[-1mm]\small Observo, escucho y pregunto.} &
\phasecard{2}{milcTurquesa}{milcGris}{Sistematización\\[-1mm]\small Pensamiento computacional.} &
\phasecard{3}{milcMagenta}{milcGris}{Praxis\\[-1mm]\small Construyo y aplico.} &
\phasecard{4}{milcVino}{milcGris}{Evaluación\\[-1mm]\small Reflexión liberadora.}
\end{tabularx}


\puente{Antes de teorizar, vas a hacer un \emph{experimento simple}: pídele a ChatGPT, Gemini o Bing que te escriba un correo a tu rector solicitando una reunión sobre el horario del recreo. Lee la respuesta. ¿Suena a ti o suena a chatbot genérico?}

\begin{titlebox}{milcVerde}
Fase 1 · Escuta
\end{titlebox}

\begin{softbox}{milcVerde}{milcGris}{Escena para escuchar}
\textbf{Actividad 1 · IDENTIFICA --- Experimento del correo IA puro} (15 min · individual con dispositivo). Abre ChatGPT, Gemini o Bing Copilot. Pídele: \emph{``Escríbeme un correo formal al rector de mi colegio solicitando reunión para discutir el horario del recreo''}. Lee la respuesta completa. Anota: (1) ¿usa fórmulas típicas de chatbot (\emph{``espero que se encuentre bien''}, \emph{``en pleno siglo XXI''})? (2) ¿el asunto es claro y específico? (3) ¿el pedido aparece pronto o se esconde? (4) ¿suena a estudiante colombiano de 16 años o a manual de protocolo?
\end{softbox}

\checkbox Le pedí a la IA que escriba un correo formal.\\
\checkbox Leí la respuesta completa.\\
\checkbox Identifiqué qué partes suenan a chatbot y cuáles a persona.

\begin{infoband}{milcVerde}
\textbf{Lo que va al cuaderno (Actividad 1 · IDENTIFICA).} Título: «Actividad 1 --- Correo IA puro». \textbf{Formato:} captura del correo IA + 4 observaciones críticas. \textbf{Sabes que terminaste cuando} ves que la IA produce borrador pero no versión final.
\end{infoband}

\puente{Ya tienes un correo IA. Ahora vas a entender la \textbf{anatomía del correo profesional}: los 5 elementos, los 5 errores típicos (saludo equivocado, asunto vago, sin pedido claro, demasiado largo, firma incompleta) y cómo editar la versión IA para que suene a ti.}

\begin{titlebox}{milcTurquesa}
Fase 2 · Sistematización con pensamiento computacional
\end{titlebox}

\begin{softbox}{milcTurquesa}{milcGris}{Cuatro pilares aplicados al tema}
Tu lectura crítica te mostró los límites de la IA. Ahora necesitas la \textbf{anatomía profesional} del correo y cómo trabajar con IA como asistente.
\end{softbox}

\small
\begin{tabularx}{\textwidth}{>{\bfseries\RaggedRight\arraybackslash}p{3.6cm}Y}
\rowcolor{milcTurquesa!90}\color{white}Pilar & \color{white}\bfseries Aplicación al tema\\
1. Descomposición & Los \textbf{5 elementos obligatorios} del correo profesional se descomponen así: (1) \textbf{Asunto claro}: máximo 5-7 palabras que digan qué espera el destinatario. \emph{Mal:} \emph{``Hola''}, \emph{``Buenas''}. \emph{Bien:} \emph{``Solicitud reunión: horario recreo''}, \emph{``Propuesta mejora cafetería con anexo''}. (2) \textbf{Saludo}: formal (\emph{``Estimado Sr. Pérez:''} con dos puntos al final), semi-formal (\emph{``Hola Carlos,''}), informal (\emph{``Hola,''}). (3) \textbf{Cuerpo}: 3-5 párrafos cortos. Párrafo 1: identidad y motivo (\emph{``Soy María, estudiante grado 10. Le escribo para X''}). Párrafos 2-3: detalles concisos. Último párrafo: pedido o acción esperada (\emph{``Solicito reunión esta semana''} o \emph{``Quedo atenta a su respuesta''}). (4) \textbf{Cierre}: \emph{``Cordialmente''}, \emph{``Atentamente''}, \emph{``Quedo atento/a''}. (5) \textbf{Firma}: nombre completo + rol + institución + correo de contacto.\\
2. Reconocimiento de patrones & Los \textbf{3 tonos básicos} según destinatario: (a) \textbf{Formal}: rector, director, autoridad, persona que no conoces. Tratamiento de \emph{usted}, saludo completo, sin emojis. (b) \textbf{Semi-formal}: docente, coordinador, alguien que conoces poco. Tratamiento de \emph{usted} pero más cálido, saludo con nombre. (c) \textbf{Cordial-informal}: compañero que conoces, profesor cercano. Tratamiento de \emph{tú} si es costumbre, saludo amistoso, ocasional emoji aceptable. La regla profesional: \textbf{ante la duda, usa el tono más formal de los 2 posibles}.\\
3. Abstracción & Los \textbf{5 errores típicos} del correo profesional: (a) \textbf{Asunto vago}: \emph{``Una cosita''}, \emph{``Pregunta''}. (b) \textbf{Demasiado largo}: 500 palabras cuando 150 bastaban. (c) \textbf{Sin pedido claro}: cierre con \emph{``espero pueda considerarlo''} sin decir qué se pide. (d) \textbf{Mezcla de registros}: empezar formal y terminar informal. (e) \textbf{Firma incompleta}: solo nombre, sin rol ni institución.\\
4. Algoritmo conceptual & El \textbf{flujo con IA como asistente} sigue 4 pasos: (1) Escribir prompt con las 5 partes (sesión 3: rol, contexto, tarea, formato, restricciones). (2) Recibir borrador de IA. (3) Editar a mano: corregir saludo si es muy formal/informal, ajustar tono colombiano, eliminar clichés chatbot, agregar pedido claro si falta, firmar con tus datos. (4) Releer en voz alta antes de enviar. La phronesis del oficio consiste en \textbf{usar IA para velocidad, conservar voz humana para profesionalismo}.\\
\end{tabularx}
\normalsize

\begin{drawbox}{milcTurquesa}{Anatomía del correo profesional}
\textbf{Asunto:} 5-7 palabras específicas. \\[1mm] \textbf{Saludo:} formal / semi-formal / cordial. \\[1mm] \textbf{Cuerpo:} 3-5 párrafos cortos + pedido claro. \\[1mm] \textbf{Cierre:} cordialmente / atentamente. \\[1mm] \textbf{Firma:} nombre + rol + institución + contacto.
\end{drawbox}

\begin{softbox}{milcMostaza}{milcGris}{Errores comunes y su corrección}
(1) Asunto vago. (2) Correo demasiado largo. (3) Sin pedido claro. (4) Mezcla de registros formal/informal. (5) Firma incompleta. (6) Aceptar borrador de IA sin editar. (7) Mantener clichés de chatbot (\emph{``espero que se encuentre bien''} usado mecánicamente).
\end{softbox}

\begin{infoband}{milcTurquesa}
\textbf{Lo que va al cuaderno (Actividad 2 · APLICA).} Título: «Actividad 2 --- Anatomía correo profesional». \textbf{Formato:} (a) 5 elementos con ejemplo; (b) 3 tonos según destinatario; (c) 7 errores con marca. \textbf{Sabes que terminaste cuando} puedes redactar correo profesional sin abrir esta guía.
\end{infoband}

\puente{Tienes el método. Toca aplicarlo: escribes 3 correos profesionales en 3 situaciones distintas (solicitar, responder, declinar), cada uno con borrador IA + versión final tuya + bitácora.}

\begin{titlebox}{milcMagenta}
Fase 3 · Praxis con pensamiento computacional
\end{titlebox}

\begin{softbox}{milcMagenta}{milcGris}{Construyendo el producto con los cuatro pilares}
\textbf{Actividad 3 · ANALIZA + EVALÚA --- 3 correos + bitácora + ficha} (40 min · individual con dispositivo). Hora de aplicar. Escribes 3 correos en 3 situaciones distintas con flujo IA-asistente, bitácora documentando cambios, ficha de 5 elementos para uso futuro.
\end{softbox}

\small
\begin{tabularx}{\textwidth}{>{\bfseries\RaggedRight\arraybackslash}p{3.6cm}Y}
\rowcolor{milcMagenta!90}\color{white}Pilar & \color{white}\bfseries Aplicación al producto\\
Descomposición & Los 3 correos a escribir representan 3 situaciones profesionales típicas: (1) \textbf{Solicitar información}: a una institución o persona pidiéndole datos, agenda, contacto. Ej. correo a la Cámara de Comercio pidiendo datos para tu estudio de mercado. (2) \textbf{Responder a un cliente o usuario}: contestar una consulta o queja con respeto y solución. Ej. respuesta a un cliente del microemprendimiento familiar. (3) \textbf{Declinar una propuesta con respeto}: rechazar amablemente algo (oferta, invitación, propuesta) sin ofender. Es la situación más difícil. Ej. declinar invitación a participar en un evento que no te conviene.\\
Patrones & El \textbf{flujo de cada correo} sigue 5 pasos: (1) Escribir prompt completo con las 5 partes. (2) Recibir borrador de IA en ChatGPT/Gemini/Bing/Claude. (3) Copiar borrador IA a tu cuaderno. (4) Editar a mano la versión final con tu voz, ajustando tono y pedido. (5) Documentar en bitácora: prompt + versión IA + versión final + razones de cambio.\\
Abstracción & La bitácora de cada correo tiene 4 columnas: (a) \textbf{Prompt usado}: el texto exacto que le diste a la IA. (b) \textbf{Versión IA}: pegada tal cual la dio. (c) \textbf{Versión final tras tu edición}: tu correo definitivo. (d) \textbf{Razones de cambio}: \emph{``cambié el saludo porque era muy formal''}, \emph{``eliminé el cliché del primer párrafo''}, \emph{``agregué pedido claro al final''}.\\
Algoritmo de armado & Algoritmo: (1) escribir 3 prompts. (2) generar 3 borradores con IA. (3) editar 3 versiones finales. (4) bitácora para cada uno. (5) ficha de 5 elementos del correo profesional ideal para uso futuro.\\
\end{tabularx}
\normalsize

\begin{drawbox}{milcMagenta}{Checklist antes de entregar}
\checkbox 3 correos en 3 situaciones distintas. \\[1mm] \checkbox Borradores IA generados con prompt profesional. \\[1mm] \checkbox 3 versiones finales editadas a mano. \\[1mm] \checkbox Bitácora con 4 columnas por cada correo. \\[1mm] \checkbox Ficha de 5 elementos para uso futuro. \\[1mm] \checkbox Cada correo con asunto claro y pedido visible.
\end{drawbox}

\begin{softbox}{milcMagenta}{milcGris}{Plantilla de trabajo}
\textbf{Correo 1 (Solicitar).} \textit{Prompt + IA + final + cambios.} \\[1mm] \textbf{Correo 2 (Responder).} \textit{Prompt + IA + final + cambios.} \\[1mm] \textbf{Correo 3 (Declinar).} \textit{Prompt + IA + final + cambios.} \\[1mm] \textbf{Ficha ideal.} \textit{5 elementos para reutilizar.}
\end{softbox}

\begin{infoband}{milcMagenta}
\textbf{Lo que va al cuaderno (Actividad 3 · ANALIZA + EVALÚA).} Título: «Actividad 3 --- 3 correos profesionales». \textbf{Formato:} (a) los 3 correos completos; (b) bitácora por cada uno; (c) ficha de 5 elementos. \textbf{Sabes que terminaste cuando} cualquiera de los 3 correos podrías enviarlo hoy a un destinatario real.
\end{infoband}

\puente{Lo que sigue resume \emph{qué} entregas hoy: 3 correos + bitácora + ficha de 5 elementos. El criterio principal: que cualquiera de los 3 podría enviarse hoy a un destinatario real.}

\begin{titlebox}{milcMostaza}
Producto de la sesión
\end{titlebox}
\textbf{3 correos profesionales + bitácora + ficha de 5 elementos}.
\begin{softbox}{milcMostaza}{milcGris}{Criterios de calidad}
(1) 3 correos en 3 situaciones distintas. (2) Asunto claro en cada uno. (3) Pedido visible en cada uno. (4) Versión final con voz propia (no chatbot). (5) Bitácora con 4 columnas. (6) Ficha reutilizable para uso futuro.
\end{softbox}

\puente{Antes de cerrar, mira el correo profesional desde las cinco dimensiones humanas. Escribir con precisión es respeto por el lector ocupado; mandar correos largos sin pedido claro es desperdicio del tiempo ajeno.}

\begin{titlebox}{milcVino}
Fase 4 · Evaluación liberadora — cinco dimensiones
\end{titlebox}

\begin{softbox}{milcVino}{milcGris}{Sentido de la evaluación}
Evaluar no es solo revisar una nota. Es reconocer qué aprendí, qué cambió en mí, qué emociones manejé, cómo me relaciono con la ciudad y la comunidad, y qué le dejaré a quien viene detrás.
\end{softbox}

\textbf{Mi reflexión liberadora en cinco dimensiones.} Completa cada frase iniciada:

\small
\begin{tabularx}{\textwidth}{>{\bfseries\RaggedRight\arraybackslash}p{3.4cm}Y>{\RaggedRight\arraybackslash}p{4.6cm}}
\rowcolor{milcVino}\color{white}Dimensión & \color{white}\bfseries Mi respuesta (frame starter) & \color{white}\bfseries Algo que descubrí\\
1. Personal & Lo que más cambió en mí fue \linead{6.5cm} & \linead{4.2cm}\\
2. Emocional & La emoción más fuerte fue \linead{2.5cm} y la regulé \linead{3cm} & \linead{4.2cm}\\
3. Ciudadana & En democracia esto importa porque \linead{6cm} & \linead{4.2cm}\\
4. Local (Cartago) & En mi barrio sería útil para \linead{6.5cm} & \linead{4.2cm}\\
5. Intergeneracional & A mi abuela/abuelo le diría \linead{6.5cm} & \linead{4.2cm}\\
\end{tabularx}
\normalsize

\begin{infoband}{milcVino}
\textbf{Para la próxima sustentación.} Vuelve a esta tabla en seis meses. Verás que las respuestas cambian — no porque lo aprendido cambie, sino porque tú cambias. La evaluación liberadora es una conversación contigo a través del tiempo.
\end{infoband}


\puente{Y para terminar, tres voces sobre la palabra justa. Dussel sobre los correos que respetan al destinatario, Séneca sobre la brevedad como virtud, Floridi sobre la comunicación profesional como ética informacional.}

\begin{titlebox}{milcGradeDark}
Triángulo de pensamiento · cierre filosófico
\end{titlebox}

\begin{softbox}{milcVino}{milcGris}{Dussel · lente del nosotros}
\textit{``El correo que respeta el tiempo del destinatario es liberador; el extenso sin pedido claro es violencia comunicativa disfrazada de cortesía.''}

Quien recibe un correo de 500 palabras sin pedido claro pierde tiempo doblemente: leyendo y descifrando qué se quiere. La filosofía de la liberación pide lo opuesto: \textbf{respetar el tiempo ajeno con claridad}. La phronesis del correo profesional consiste en \textbf{decir lo justo con respeto, ni más ni menos}.

\textbf{¿Mis correos respetan el tiempo del lector, o lo desperdicio con palabras innecesarias?}
\end{softbox}
\linead{\linewidth}\\[1mm]
\linead{\linewidth}

\begin{softbox}{milcMostaza}{milcGris}{Estoicismo · Séneca · cuidado interior}
\textit{``La brevedad es virtud profesional; la extensión innecesaria es vicio del que escribe para impresionar.''}

Es tentador escribir correos largos creyendo que parecen más profesionales. La sabiduría estoica recomienda lo opuesto: \emph{breve y preciso vale más que largo y elegante}. El telegrafista cobraba por palabra; el correo profesional moderno cobra atención. Cada palabra de más cuesta atención del lector. La phronesis del oficio honra la brevedad útil.

\textbf{¿Estoy escribiendo con la brevedad necesaria, o agrego palabras para parecer más profesional?}
\end{softbox}
\linead{\linewidth}\\[1mm]
\linead{\linewidth}

\begin{softbox}{milcTurquesa}{milcGris}{Floridi · lente de la infoesfera}
\textit{``La comunicación profesional con IA bien usada es ética informacional contemporánea; depender de IA sin editar es renunciar al criterio humano.''}

En la era de la IA, muchos correos se envían directamente como salen de ChatGPT. La ética profesional pide lo contrario: usar la IA como asistente para velocidad, conservar tono humano para criterio. Tus 3 correos de hoy te entrenan en esa práctica que rige toda comunicación profesional contemporánea: \textbf{IA como herramienta, humano como autor}.

\textbf{¿Mis correos suenan a mí con apoyo de IA, o suenan a IA pura sin mi criterio?}
\end{softbox}
\linead{\linewidth}\\[1mm]
\linead{\linewidth}

\begin{titlebox}{milcVino}
Compromiso final
\end{titlebox}

\small
\begin{tabularx}{\textwidth}{>{\RaggedRight\arraybackslash}p{3.0cm}YYY}
\rowcolor{milcVino}\color{white}\bfseries Criterio & \color{white}\bfseries Lo logré & \color{white}\bfseries En proceso & \color{white}\bfseries Necesito apoyo\\
Comprensión del tema & Explico ideas centrales con mis palabras. & Reconozco ideas pero falta precisión. & Necesito repasar conceptos base.\\
Pensamiento computacional & Apliqué los 4 pilares al diseñar mi producto. & Apliqué algunos pilares. & Necesito releer la fase 2.\\
Aplicación y producto & Mi producto cumple los criterios y se puede revisar. & Producto funcional con ajustes pendientes. & Aún en construcción.\\
Reflexión 5 dimensiones & Respondí cada dimensión con honestidad. & Respondí algunas con detalle. & Mis respuestas son muy generales.\\
Triángulo de pensamiento & Conecté las 3 voces con mi trabajo real. & Conecté al menos una. & No relaciono las voces con mi proceso.\\
\end{tabularx}
\normalsize

\textbf{Hoy entendí que:}\\[1mm]
\linead{\linewidth}\\[1mm]
\linead{\linewidth}

\textbf{Mi compromiso para la próxima sesión es:}\\[1mm]
\linead{\linewidth}\\[1mm]
\linead{\linewidth}

\begin{softbox}{milcMostaza}{milcGris}{Cita de despedida · Marco Aurelio}
\textit{``El obstáculo en el camino se vuelve el camino. Lo que estorba al camino se vuelve camino.''}\\[1mm]
Cada error de hoy, bien observado, se convierte en pieza del camino que pisarás mañana.
\end{softbox}

\small
\begin{tabularx}{\textwidth}{>{\bfseries}p{3.6cm}Y}
\rowcolor{milcGradeDark!12}Nombre completo: & \linead{10cm}\\
Fecha: & \linead{4cm} \quad Curso/grupo: \linead{4cm}\\
Mi firma: & \linead{9cm}\\
\end{tabularx}
\normalsize

\begin{infoband}{milcGradeDark}
\textbf{Para la próxima sesión.} Llega con tus 3 correos profesionales y la bitácora. En la sesión 9 vas a integrar todo el periodo en un \textbf{mini-proyecto: informe técnico con datos reales y IA documentada}. Es la cima del P2 antes de la sustentación de S10.
\end{infoband}

\end{document}
